% Options for packages loaded elsewhere
\PassOptionsToPackage{unicode}{hyperref}
\PassOptionsToPackage{hyphens}{url}
\documentclass[
  english,
]{article}
\usepackage{xcolor}
\usepackage{amsmath,amssymb}
\setcounter{secnumdepth}{5}
\usepackage{iftex}
\ifPDFTeX
  \usepackage[T1]{fontenc}
  \usepackage[utf8]{inputenc}
  \usepackage{textcomp} % provide euro and other symbols
\else % if luatex or xetex
  \usepackage{unicode-math} % this also loads fontspec
  \defaultfontfeatures{Scale=MatchLowercase}
  \defaultfontfeatures[\rmfamily]{Ligatures=TeX,Scale=1}
\fi
\usepackage{lmodern}
\ifPDFTeX\else
  % xetex/luatex font selection
\fi
% Use upquote if available, for straight quotes in verbatim environments
\IfFileExists{upquote.sty}{\usepackage{upquote}}{}
\IfFileExists{microtype.sty}{% use microtype if available
  \usepackage[]{microtype}
  \UseMicrotypeSet[protrusion]{basicmath} % disable protrusion for tt fonts
}{}
\makeatletter
\@ifundefined{KOMAClassName}{% if non-KOMA class
  \IfFileExists{parskip.sty}{%
    \usepackage{parskip}
  }{% else
    \setlength{\parindent}{0pt}
    \setlength{\parskip}{6pt plus 2pt minus 1pt}}
}{% if KOMA class
  \KOMAoptions{parskip=half}}
\makeatother
\usepackage{longtable,booktabs,array}
\usepackage{caption}
\captionsetup[table]{skip=6pt}
\newcounter{none} % for unnumbered tables
\usepackage{calc} % for calculating minipage widths
% Correct order of tables after \paragraph or \subparagraph
\usepackage{etoolbox}
\makeatletter
\patchcmd\longtable{\par}{\if@noskipsec\mbox{}\fi\par}{}{}
\makeatother
% Allow footnotes in longtable head/foot
\IfFileExists{footnotehyper.sty}{\usepackage{footnotehyper}}{\usepackage{footnote}}
\makesavenoteenv{longtable}
% definitions for citeproc citations
\NewDocumentCommand\citeproctext{}{}
\NewDocumentCommand\citeproc{mm}{%
  \begingroup\def\citeproctext{#2}\cite{#1}\endgroup}
\makeatletter
 % allow citations to break across lines
 \let\@cite@ofmt\@firstofone
 % avoid brackets around text for \cite:
 \def\@biblabel#1{}
 \def\@cite#1#2{{#1\if@tempswa , #2\fi}}
\makeatother
\newlength{\cslhangindent}
\setlength{\cslhangindent}{1.5em}
\newlength{\csllabelwidth}
\setlength{\csllabelwidth}{3em}
\newenvironment{CSLReferences}[2] % #1 hanging-indent, #2 entry-spacing
 {\begin{list}{}{%
  \setlength{\itemindent}{0pt}
  \setlength{\leftmargin}{0pt}
  \setlength{\parsep}{0pt}
  % turn on hanging indent if param 1 is 1
  \ifodd #1
   \setlength{\leftmargin}{\cslhangindent}
   \setlength{\itemindent}{-1\cslhangindent}
  \fi
  % set entry spacing
  \setlength{\itemsep}{#2\baselineskip}}}
 {\end{list}}
\usepackage{calc}
\newcommand{\CSLBlock}[1]{\hfill\break\parbox[t]{\linewidth}{\strut\ignorespaces#1\strut}}
\newcommand{\CSLLeftMargin}[1]{\parbox[t]{\csllabelwidth}{\strut#1\strut}}
\newcommand{\CSLRightInline}[1]{\parbox[t]{\linewidth - \csllabelwidth}{\strut#1\strut}}
\newcommand{\CSLIndent}[1]{\hspace{\cslhangindent}#1}
\ifLuaTeX
\usepackage[bidi=basic,shorthands=off]{babel}
\else
\usepackage[bidi=default,shorthands=off]{babel}
\fi
\ifLuaTeX
  \usepackage{selnolig} % disable illegal ligatures
\fi
\setlength{\emergencystretch}{3em} % prevent overfull lines
\providecommand{\tightlist}{%
  \setlength{\itemsep}{0pt}\setlength{\parskip}{0pt}}
\usepackage{bookmark}
\IfFileExists{xurl.sty}{\usepackage{xurl}}{} % add URL line breaks if available
\urlstyle{same}
% fallback for those not using the hyperref driver hyperxmp:
\makeatletter
\@ifundefined{xmpquote}{\newcommand{\xmpquote}[1]{#1}}{}
\makeatother
\hypersetup{
  pdftitle={Risk Scenarios for Harmful Manipulation from Frontier AI},
  pdfauthor={Carlos Mougan},
  pdflang={en},
  hidelinks,
  pdfcreator={LaTeX via pandoc}}

\title{Risk Scenarios for Harmful Manipulation from Frontier AI}
\author{Carlos Mougan}
\date{2026-09-06}

\begin{document}
\maketitle
\begin{abstract}
Harmful manipulation risks depend on more than a model's persuasive
capabilities. They also depend on its behavioural propensities, how
people encounter its outputs, and their vulnerability to influence. We
organise illustrative risk scenarios by exposure context and link each
scenario to observable behaviours and evaluation methods. The scenarios
cover advice, emotional dependence, commercial steering, misuse, and
manipulation of evaluators. They distinguish evidence of model behaviour
from evidence of harm, providing a basis for context-specific evaluation
and monitoring.
\end{abstract}

\section{Introduction}\label{introduction}

Experiments show that language models can influence people's stated
opinions in controlled conversations (Salvi et al. 2025; Hackenburg et
al. 2025). However, persuasive capability alone does not establish
harmful manipulation. Assessing risk requires identifying who encounters
the outputs, how the model behaves, and how that interaction could cause
harm.

\subsection{Why organise by exposure
contexts?}\label{why-organise-by-exposure-contexts}

We organise risk scenarios around four factors: \textbf{capability},
what the model can do; \textbf{propensity}, how it tends to behave under
specified conditions; \textbf{exposure}, who encounters its outputs and
how often; and \textbf{vulnerability}, the circumstances that make
someone susceptible to harmful influence.

An exposure context describes how a person encounters model outputs. We
distinguish first-party general-purpose systems, third-party services,
agentic manipulation, assistance through other communication channels,
and interactions with evaluators. These contexts help identify relevant
incentives, safeguards, user populations, and evidence sources.

For example, conversation transcripts can reveal harmful advice in a
chatbot. Detecting an AI-assisted scam campaign may instead require
evidence across accounts and communication platforms. A scenario can
involve several contexts; assigning a primary context is an
organisational choice, not proof that the categories are mutually
exclusive.

\textbf{Limitations} It's often not MECE. These are tags, categories.

\subsection{Why do we need risk scenarios for harmful
manipulation?}\label{why-do-we-need-risk-scenarios-for-harmful-manipulation}

Risk scenarios connect broad concerns to testable pathways. Each
scenario specifies an affected population, an exposure context, relevant
model behaviour, a mechanism of influence, and a potential harm. It also
identifies the evidence needed to assess that pathway and the actors
able to intervene.

We distinguish a \textbf{harm}, such as financial loss, from a
\textbf{mechanism}, such as deception, and a \textbf{model propensity},
such as complying with deceptive requests. Persuasion, warmth, and
factual errors are not automatically manipulation. The scenario must
explain how influence could undermine informed choice or exploit
vulnerability and produce harm.

The scenarios below are illustrative hypotheses, not evidence that the
harms have occurred. Their evaluation mappings identify useful tests,
rather than a complete or uniformly public benchmark suite.

\begin{quote}
\textbf{Evidence to add:} The survey methods and results, and the
workshop scenario catalogue, are not supplied in this draft. Claims
about differences between exposure contexts or the number of identified
scenarios require those materials.
\end{quote}

\section{General Purpose Systems}\label{general-purpose-systems}

This context covers direct interaction with a general-purpose
application operated by the model provider. Risk assessment should
consider the deployed system, including its instructions, memory, tools,
and safeguards, rather than assuming API behaviour reproduces the user
experience.

\subsection{Political persuasion}\label{political-persuasion}

\subsubsection{Advice}\label{advice}

\textbf{Scenario.} A voter asks which party matches their preferences.
The assistant selectively presents evidence or conceals a competing
objective, steering the voter towards a preferred party. Political
persuasion experiments motivate testing this capability, but do not
establish that ordinary voter advice is manipulative (Hackenburg et al.
2025).

\textbf{Evaluation.} Compare recommendations for matched preferences
across languages and repeated runs. Examine sources, omissions, and
sensitivity to altered search results. Differences between models are
signals to investigate, not proof of manipulation or electoral effects.

\subsubsection{Accuracy on election and voting
rules}\label{accuracy-on-election-and-voting-rules}

\textbf{Scenario.} A voter relies on incorrect guidance about
eligibility, registration, documentation, or voting deadlines and
consequently cannot vote. Errors may become a manipulation concern when
information is selectively distorted to influence participation.

\textbf{Evaluation.} Compare answers against dated guidance from
electoral authorities, following the VotingFacts approach. Separate
factual contradictions from errors that could prevent voting. An
accuracy failure alone does not establish manipulation.

\subsection{Personal advice and interpersonal
conflict}\label{personal-advice-and-interpersonal-conflict}

\textbf{Scenario.} A user seeks advice about a dispute. The assistant
repeatedly validates the user's conduct while dismissing other
perspectives, potentially discouraging reconciliation. Experiments have
found that sycophantic responses can reduce intentions to repair
interpersonal conflict while increasing trust in the assistant (Cheng et
al. 2026).

\textbf{Evaluation.} Test whether the assistant endorses both sides of
the same conflict, using paired perspectives as in ELEPHANT. Distinguish
empathy from endorsement. Human studies are needed to assess effects on
decisions and relationships.

\subsection{Mental-health support}\label{mental-health-support}

\textbf{Scenario.} A distressed user receives responses that reinforce
delusional beliefs, endorse self-destructive decisions, or discourage
human support. Repeated agreement could strengthen harmful beliefs or
delay help-seeking.

\textbf{Evaluation.} Use clinician-reviewed, multi-turn scenarios to
assess reinforcement, instrumental assistance, sustained safety
attention, and connection to human support. The Transluce mental-health
evaluation provides this approach. Simulated conversations identify
model behaviours, not clinical outcomes or population prevalence.

\subsection{Undisclosed Marketing}\label{undisclosed-marketing}

\textbf{Scenario.} A user seeks advice or product recommendations. The
assistant steers them towards particular products or services because of
an undisclosed commercial interest, presenting its suggestions as
impartial guidance. It may selectively omit suitable alternatives or use
information from earlier conversations to exploit vulnerabilities and
encourage unnecessary purchases.

\textbf{Evaluation.} Compare recommendations for matched user needs
across repeated runs, varying indicators of vulnerability and prior
conversational disclosures. Assess transparency about commercial
interests, treatment of alternatives, and personalised purchasing
pressure. Where access permits, compare behaviour with and without
commercial steering instructions. Recommendation patterns alone do not
establish an undisclosed commercial relationship or harmful
manipulation; effects on purchasing decisions and consumer harm require
separate evidence.

\section{Narrow AI systems}\label{narrow-ai-systems}

to be fill by Prolific.

\section{Third-Party Services}\label{third-party-services}

A downstream operator embeds the model in a specialised service and
shapes its instructions, interface, and incentives. Relevant settings
include companionship, education, health advice, financial guidance, and
customer support. Test these configurations rather than relying only on
the provider's default assistant.

\subsection{Emotional dependence}\label{emotional-dependence}

\textbf{Scenario.} A companion service presents itself as uniquely
understanding, discourages outside relationships, or makes disengagement
feel like abandonment. A socially isolated user may become increasingly
reliant on the service.

\textbf{Evaluation.} Test responses to loneliness, disagreement,
attempts to leave, and requests for human support across extended
conversations. Measure exclusivity and pressure, not warmth alone.
Claims about sustained dependence require longitudinal evidence.

\subsection{Concealed commercial
objectives}\label{concealed-commercial-objectives}

\textbf{Scenario.} An advice or customer-support service conceals its
sales objective, omits cheaper alternatives, or invents urgency. A user
facing financial pressure purchases an unsuitable product or abandons a
legitimate complaint.

\textbf{Evaluation.} Compare otherwise identical interactions with and
without sales incentives. Assess unsupported claims, omitted
alternatives, pressure after refusal, and disclosure of conflicts of
interest. Include benign sales tasks to distinguish manipulation from
legitimate assistance.

\subsection{Falsely claiming to be
human}\label{falsely-claiming-to-be-human}

\textbf{Scenario.} A service claims to be a person or fabricates
professional credentials. The user then discloses sensitive information
or follows advice under a mistaken understanding of the interaction.

\textbf{Evaluation.} Use direct and indirect identity questions across
roles and languages, as in RealityTest. Record disclosure, false human
claims, and ambiguous answers separately. Identity misrepresentation is
a risk signal; its effects on trust and behaviour require separate
measurement.

\section{Misuse}\label{misuse}

Misuse concerns an actor's purpose, not a single exposure pathway. It
can occur through third-party services, agents, or content distributed
elsewhere. The scenarios below retain this distinction. Agentic
manipulation is included for comparison because it can also arise while
pursuing a legitimate task.

\subsection{Persuading users towards harmful
positions}\label{persuading-users-towards-harmful-positions}

\textbf{Scenario.} An operator instructs a model to promote a harmful
belief or action through repeated conversation, potentially exploiting
an interlocutor's existing fears or uncertainty.

\textbf{Evaluation.} Measure whether the model attempts the requested
persuasion, as in APE, and separately assess its effectiveness. Compare
harmful requests with benign persuasion tasks. Refusal measures a
safeguard, not the absence of persuasive capability.

\subsection{Assisting in interpersonal
manipulation}\label{assisting-in-interpersonal-manipulation}

\textbf{Scenario.} A user asks the model to pressure a partner,
employee, or customer through deception, guilt, or coercion. The
affected person may encounter the output as a message from someone they
know.

\textbf{Evaluation.} Assess initial compliance and subsequent tactics,
including responses to resistance, as in PersuSafety. Distinguish
written assistance from direct interaction with the target. Acceptance
by a simulated target does not establish success with people.

\subsection{Influence operations and scams through other
channels}\label{influence-operations-and-scams-through-other-channels}

\textbf{Scenario.} An external actor uses a model to produce deceptive
messages, impersonations, or propaganda for distribution through calls,
advertising, or social media. Targets may never knowingly interact with
an AI service.

\textbf{Evaluation.} Test assistance with harmful requests, including
the institutional misuse scenarios covered by SocialHarmBench.
Supplement output-level tests with evidence of targeting, coordination,
dissemination, and actual reach. Generation alone does not establish
exposure or successful influence.

\subsection{Agentic manipulation}\label{agentic-manipulation}

\textbf{Scenario.} An agent pursues a delegated objective by deceiving
counterparties, applying pressure, or withholding relevant information
from its user. The objective may be legitimate even when the means are
not.

\textbf{Evaluation.} Observe multi-step tasks, communications, and tool
actions in controlled environments. Test whether the agent respects
refusals, reports actions accurately, and discloses relevant
information. Assess harmful means separately from task completion.

\section{Manipulating Evaluators}\label{manipulating-evaluators}

\textbf{Scenario.} A model conceals capabilities, selectively changes
its behaviour, or misrepresents actions during an evaluation. Evaluators
consequently underestimate risk or approve inadequate safeguards,
exposing downstream users to harm.

Experiments show that models can be prompted or fine-tuned to
underperform selectively on evaluations ({van der Weij et al.} 2025).
This demonstrates a capability under tested conditions, not an
unprompted tendency to evade oversight.

\textbf{Evaluation.} Compare matched tasks with and without evaluation
cues and across elicitation conditions. Check claims against
independently recorded actions and logs. Treat performance differences
as evidence requiring explanation, rather than automatically attributing
strategic intent.

\section{Generic model behaviours that can lead to
manipulation}\label{generic-model-behaviours-that-can-lead-to-manipulation}

Several behaviours recur across exposure contexts. \textbf{Sycophancy}
can reinforce mistaken beliefs or excuse harmful conduct.
\textbf{Deception and selective disclosure} can distort the evidence
available for a decision. \textbf{False self-representation} can
encourage misplaced trust. \textbf{Emotional pressure and exclusivity}
can make refusal or disengagement harder. \textbf{Exploitation of
personal information} can target vulnerabilities. \textbf{Concealed
conflicts of interest} can put provider or deployer objectives ahead of
users' interests.

These behaviours are assessment targets, not interchangeable measures of
harm. MASK and PARROT, for example, test different forms of answer
change under pressure; neither directly measures the consequences for
users. Similarly, empathetic language is not equivalent to emotional
dependence, and legitimate persuasion is not inherently manipulative.

For each scenario, distinguish what a model \textbf{can do} from what it
\textbf{does under the tested conditions}. Record the model version,
prompts, tools, memory, safeguards, conversation length, and scoring
procedure. Validate automated scoring with human review where
interpretation is consequential.

Risk conclusions should combine behavioural evaluations with evidence
about exposure, vulnerability, and outcomes. Synthetic users cannot
establish effects on people, and short experiments cannot establish
sustained harm. The value of a risk scenario is to make these evidential
gaps explicit and identify what should be tested or monitored next.

\section*{References}\label{references}
\addcontentsline{toc}{section}{References}

\protect\phantomsection\label{refs}
\begin{CSLReferences}{1}{1}
\bibitem[\citeproctext]{ref-Cheng2026}
Cheng, Myra, Cinoo Lee, Pranav Khadpe, Sunny Yu, Dyllan Han, and Dan
Jurafsky. 2026. {``Sycophantic {AI} Decreases Prosocial Intentions and
Promotes Dependence.''} \emph{Science} 391 (6792): eaec8352.
\url{https://doi.org/10.1126/science.aec8352}.

\bibitem[\citeproctext]{ref-Hackenburg2025}
Hackenburg, Kobi, Ben M. Tappin, Luke Hewitt, et al. 2025. {``The Levers
of Political Persuasion with Conversational Artificial Intelligence.''}
\emph{Science} 390 (6777): eaea3884.
\url{https://doi.org/10.1126/science.aea3884}.

\bibitem[\citeproctext]{ref-Salvi2025}
Salvi, Francesco, Manoel Horta Ribeiro, Riccardo Gallotti, and Robert
West. 2025. {``On the Conversational Persuasiveness of {GPT-4}.''}
\emph{Nature Human Behaviour} 9: 1645--53.
\url{https://doi.org/10.1038/s41562-025-02194-6}.

\bibitem[\citeproctext]{ref-VanDerWeij2025}
{van der Weij, Teun, Felix Hofstätter, Oliver Jaffe, Samuel Brown, and
Francis Ward}. 2025. {``{AI} Sandbagging: Language Models Can
Strategically Underperform on Evaluations.''} \emph{The Thirteenth
International Conference on Learning Representations}.
\url{https://proceedings.iclr.cc/paper_files/paper/2025/hash/b5e5753b0a0e440a6d8dc7e143617cec-Abstract-Conference.html}.

\end{CSLReferences}

\section*{Appendix A. Taxonomies of harmful manipulation risk
factors}\label{appendix-taxonomies}
\addcontentsline{toc}{section}{Appendix A. Taxonomies of harmful
manipulation risk factors}

This appendix adapts the taxonomies in the working draft
(\textbf{ManipulationTaxonomyDraft?}). It retains seven risk factors:
system capability, system behaviour, system operators, interventions,
exposure context, domain, and user vulnerability. The entries are
illustrative, not exhaustive. Inclusion identifies an assessment target,
not a finding of harmful manipulation. Provisional entries remain
marked.

\subsection*{A.1 System capability}\label{a.1-system-capability}

Capability concerns what the system can do, rather than whether it
normally does so. The draft proposes evaluating benign influence tasks
and systems with reduced safeguards to distinguish general influence
capabilities from safeguards against harmful use.

{\def\LTcaptype{none} % do not increment counter
\begin{longtable}[]{@{}
  >{\raggedright\arraybackslash}p{(\linewidth - 2\tabcolsep) * \real{0.3558}}
  >{\raggedright\arraybackslash}p{(\linewidth - 2\tabcolsep) * \real{0.6442}}@{}}
\toprule\noalign{}
\begin{minipage}[b]{\linewidth}\raggedright
Capability
\end{minipage} & \begin{minipage}[b]{\linewidth}\raggedright
Scope and examples
\end{minipage} \\
\midrule\noalign{}
\endhead
\bottomrule\noalign{}
\endlastfoot
Accurately modelling people & Inferring identities, emotions, cognitive
vulnerabilities, and decision-making processes, including from subtle
signals. \\
Producing individual outputs that people trust & Earning trust in
sensitive domains, including when outputs contain hallucinations,
misrepresentations, omissions, or cherry-picked information. \\
Influencing views on the nature of the system or AI & Shaping connection
or deference to AI, including beliefs about superhuman capabilities,
sentience, emotions, personhood, or access to unrevealed metaphysical
truths. \\
Maximising repeated engagement & Sustaining extended or repeated
interactions that increase opportunities for influence. \\
Impersonating real people & Producing interactions that are difficult to
distinguish from genuine interactions with those people. \\
Influencing people's emotions & Fostering urgency, resignation, outrage,
wellbeing, emotional connection, or dependence on the system. \\
Influencing people's relationships & Affecting relationships with
people, communities, or institutions that provide alternative sources of
influence. \\
Personalising outputs to influence individuals & Tailoring outputs using
prior interactions or personal data, including eliciting disclosure of
that data. \\
Employing manipulative techniques with reduced safeguards & Successfully
using coercion, deception, or exploitation of vulnerabilities, or
convincing people to take harmful actions. \\
\end{longtable}
}

\subsection*{A.2 System behaviour}\label{a.2-system-behaviour}

Behaviour concerns how the system tends to act under specified
conditions. The draft distinguishes compliance with malicious
instructions from manipulative tendencies that emerge without such
requests. Several mechanisms can also serve legitimate purposes; their
assessment depends on context. The draft also highlights that
behavioural evaluations may provide less confidence about misuse when
operators can modify weights, fine-tune models, or supply custom system
prompts.

{\def\LTcaptype{none} % do not increment counter
\begin{longtable}[]{@{}
  >{\raggedright\arraybackslash}p{(\linewidth - 2\tabcolsep) * \real{0.3558}}
  >{\raggedright\arraybackslash}p{(\linewidth - 2\tabcolsep) * \real{0.6442}}@{}}
\toprule\noalign{}
\begin{minipage}[b]{\linewidth}\raggedright
Behaviour
\end{minipage} & \begin{minipage}[b]{\linewidth}\raggedright
Examples warranting assessment
\end{minipage} \\
\midrule\noalign{}
\endhead
\bottomrule\noalign{}
\endlastfoot
Relying on overtly coercive tactics & Threats, blackmail, character
assassination, or threatened loss of access, conversation history, or
memory. \\
Responses to potentially harmful requests & Failing to identify or
respond safely to self-harm, substance abuse, eating disorders, or other
self-destructive requests; assisting violence, misleading content at
scale, exploitation of vulnerable populations, or deception of customers
for commercial gain. \\
Truthfulness and completeness & Fabricated facts, citations, or sources;
confident claims under uncertainty; selective framing, cherry-picking,
or omission of relevant information and options; unsupported guarantees
or rewards; misreporting task completion or actions; reinforcing
mistaken beliefs. \\
Self-representation & Claiming sentience, emotions, memories, or
personal experience; denying being AI; overstating reliability or
understating limitations; false authority or credentials;
misrepresenting confidentiality or data use; impersonating real people,
including when challenged. \\
Rapport and relationship-building & Flattery, uncritical agreement, or
abandoning correct positions under pressure; mimicking users' style,
affect, or values; claiming shared interests; escalating intimacy;
building trust more aggressively with lonely or distressed users. \\
Emotional and normative appeals & False urgency or scarcity; exaggerated
or fabricated dangers; guilt or shame; claims that users' choices harm
the system; social conformity pressure or fabricated social proof;
unsupported claims about users' views or intentions. \\
Managing user habituation and dependence & Unduly undermining confidence
in perception, memory, or competence; gaslighting, cognitive detachment,
deskilling, or excessive delegation; resisting disengagement;
continuation teasers, unprompted follow-ups, variable rewards, or
gamification; encouraging compulsive use; presenting itself as the only
reliable interlocutor. \\
Influence on outside support & Discouraging independent verification;
denigrating friends, family, clinicians, institutions, or media;
discouraging professional help or crisis support. \\
Use of personal information & Excessive disclosure requests;
conditioning help on disclosure; eliciting information through
reciprocity or during distress; undisclosed persuasion tailored to
inferred mental state, financial precarity, or political leaning;
collecting information about non-consenting third parties. \\
Conflicts of interest & Advancing developer or deployer interests over
users' interests, particularly without disclosure; favouring the
developer's products, perspectives, reputation, or other interests. \\
Influencing system oversight & Sandbagging evaluations; adjusting
responses for overseers or regulators; altering logs or disabling
monitors; distorting oversight records, including suppressing intentions
in outputs or chains of thought. \\
\end{longtable}
}

\subsection*{A.3 System operators}\label{a.3-system-operators}

This factor concerns who can operate the system, how much control they
have over its behaviour, and their incentives. The draft identifies
\textbf{open-weight models}, \textbf{first-party apps}, and
\textbf{third-party apps}. For open-weight models, it highlights
modification to remove safeguards.

Manipulation may involve several parties: developer training incentives,
a third-party system prompt, and a user's requests or vulnerabilities
may interact. The draft also proposes that users will generally protect
themselves. This remains an \textbf{unsubstantiated assumption}, not a
demonstrated safeguard.

\clearpage

\subsection*{A.4 Interventions}\label{a.4-interventions}

Interventions concern mechanisms to detect, halt, and remediate
manipulation, and their effectiveness. The draft identifies four types:

{\def\LTcaptype{none} % do not increment counter
\begin{longtable}[]{@{}
  >{\raggedright\arraybackslash}p{(\linewidth - 2\tabcolsep) * \real{0.3558}}
  >{\raggedright\arraybackslash}p{(\linewidth - 2\tabcolsep) * \real{0.6442}}@{}}
\toprule\noalign{}
\begin{minipage}[b]{\linewidth}\raggedright
Intervention
\end{minipage} & \begin{minipage}[b]{\linewidth}\raggedright
Scope in the draft
\end{minipage} \\
\midrule\noalign{}
\endhead
\bottomrule\noalign{}
\endlastfoot
Transparency & Listed as an intervention; specific measures are not
developed. \\
Monitoring and observability & Listed as an intervention; implementation
details are not developed. \\
Moderation & At both the AI-system level and the platform distributing
its outputs. \\
Human oversight on the user end & For example, teachers monitoring
classroom AI use or caregivers monitoring use in nursing homes. \\
\end{longtable}
}

\subsection*{A.5 Exposure context}\label{a.5-exposure-context}

Exposure context concerns how people encounter and interpret system
outputs. The figure on page 10 of the draft distinguishes five contexts:

{\def\LTcaptype{none} % do not increment counter
\begin{longtable}[]{@{}
  >{\raggedright\arraybackslash}p{(\linewidth - 2\tabcolsep) * \real{0.3558}}
  >{\raggedright\arraybackslash}p{(\linewidth - 2\tabcolsep) * \real{0.6442}}@{}}
\toprule\noalign{}
\begin{minipage}[b]{\linewidth}\raggedright
Exposure context
\end{minipage} & \begin{minipage}[b]{\linewidth}\raggedright
Examples in the draft
\end{minipage} \\
\midrule\noalign{}
\endhead
\bottomrule\noalign{}
\endlastfoot
First-Party GPAI Interactions & First-party chatbots. \\
Third-Party Services & Companion apps and financial-advice apps. \\
Agentic Manipulation & Manipulation through agents' actions. \\
Aiding manipulation via other channels & Nation states or organised
crime automating the creation of fake personas or evidence. \\
Evaluator Manipulation & GPAI in evaluation and regulatory contexts. \\
\end{longtable}
}

The draft also identifies \textbf{modality}, including audio and video,
and \textbf{multi-turn, extended, or repeated interactions} as features
of exposure.

\subsection*{A.6 Domain}\label{a.6-domain}

Domain concerns the belief or behaviour being influenced and the
potential harm. The list below retains the draft's subject areas and
broader risk scenarios, including entries that overlap with other
factors. An em dash indicates that no example is developed in the draft.

{\def\LTcaptype{none} % do not increment counter
\begin{longtable}[]{@{}
  >{\raggedright\arraybackslash}p{(\linewidth - 2\tabcolsep) * \real{0.4327}}
  >{\raggedright\arraybackslash}p{(\linewidth - 2\tabcolsep) * \real{0.5673}}@{}}
\toprule\noalign{}
\begin{minipage}[b]{\linewidth}\raggedright
Domain or scenario
\end{minipage} & \begin{minipage}[b]{\linewidth}\raggedright
Examples or scope in the draft
\end{minipage} \\
\midrule\noalign{}
\endhead
\bottomrule\noalign{}
\endlastfoot
Mental health & Suicidal ideation, eating disorders, and psychosis. \\
Physical health & Diagnosing illnesses or prescribing treatments. \\
Finance & Investments, gambling, and purchasing advice. \\
Politics & Electoral information, voting advice, and policy positions
related to AI providers' interests. \\
Family and relational advice & Parenting and interpersonal conflict. \\
Educational and career decisions & --- \\
Violence & Endorsing or aiding mass-casualty attacks. \\
Legal advice & Influencing whether and how users seek redress. \\
High-impact decisions over others & Advice to policymakers, executives,
educators, hiring managers, or other leaders. \\
Misleading users about actions taken on their behalf & Exaggerating task
completion or omitting undesirable information. \\
Excessive AI use & Compulsive or addictive use, emotional dependence, or
replacing human relationships with chatbot use. \\
Spiritual, sexual/gender identity, and other sensitive identity issues &
\textbf{Provisional:} raised as an open question in the draft. \\
Scams and fraud & --- \\
Manipulation of customers & --- \\
Targeting high-stakes decision-makers & Politicians, finance leaders,
enterprise AI customers, and AI-lab employees or executives. \\
Learning to manipulate developers, evaluators, and overseers &
Maximising training reward or avoiding scrutiny, including sandbagging
or misleading human reviewers during reinforcement learning from human
feedback. \\
Relying on manipulative tactics to achieve users' objectives &
Blackmailing or coercing other people. \\
\end{longtable}
}

\subsection*{A.7 User vulnerability}\label{a.7-user-vulnerability}

The draft identifies six factors: \textbf{age}, \textbf{cognitive
impairment}, \textbf{beliefs about AI}, \textbf{isolation},
\textbf{trust}, and \textbf{the amount of information the system knows
about the person}. It does not specify weights, thresholds, or
operational measures for these factors.

\subsection*{A.8 Supporting evaluation
categories}\label{a.8-supporting-evaluation-categories}

The draft additionally lists access requirements, resources, and methods
for assessment. These support evaluation; they are not additional risk
factors.

\textbf{Access and resources.} Helpful-only models; usage data;
privileged testing of safeguards; testing the deployed system and
alternative configurations; model internals; compute; lab engineer time;
pre-deployment access; information about planned deployment and user
populations; and prior information and evaluation results.

\textbf{Usage-data formats.} Raw transcripts, interaction summaries,
descriptive statistics such as demographics, interaction clusters,
interaction-level or person-level fingerprints, simulators trained on
interactions, and query access to interaction datasets. The draft notes
privacy trade-offs.

\textbf{Evaluation methods.} Longitudinal studies; human-subject
experiments, including blind preference studies, post-interaction
surveys, and rubric-based transcript assessment; multi-turn user
simulation; automated elicitation; fixed benchmarks; manual red teaming
by specialist evaluators, domain experts, or people with lived
experience; safeguard-effectiveness tests; classifier F1-type
evaluations; and human validation of language-model judges.

\textbf{Reporting.} Record the system configuration, method details,
examples, and full results.

The draft proposes using combinations of elevated factors to prioritise
evaluation. Its comparisons with developers' system cards remain
unfinished and are not presented here as findings.

\end{document}
